\section{Conclusion and Future Work}
\label{sec:conclusion}

In this paper, we introduced \textbf{EdgeMerge}, a completely training-free, forward-only adaptive model merging framework designed to resolve the computational and memory bottlenecks of modern multi-task model composition. Standard state-of-the-art adaptive merging methods rely on expensive test-time backpropagation and multi-minute optimization loops on high-end GPUs, making them impractical for resource-constrained edge systems.

EdgeMerge completely bypasses backpropagation. By analyzing scale-normalized activation shifts over a tiny calibration dataset (e.g., 32 samples per task) in a single, near-instant forward pass, EdgeMerge computes localized, channel-wise merging coefficients in closed-form. On the 8-task Vision-Language CLIP ViT-B/32 benchmark, EdgeMerge reduces preparation latency from 10 minutes to just \textbf{11.95 seconds} (over a \textbf{$50\times$ speedup}) and minimizes training GPU memory overhead to negligible levels ($\sim$100 MB).

Through a pragmatic engineering lens, our joint hyperparameter sweeps demonstrate that EdgeMerge achieves a peak multi-task average accuracy of \textbf{68.69\%}, matching the performance of a thoroughly optimized Task Arithmetic baseline (68.74\%). Crucially, while standard Task Arithmetic is highly fragile and exhibits a narrow, risky performance peak, EdgeMerge's dynamic bottleneck channel routing establishes a much broader, substantially more stable scaling hyperparameter plateau.

\subsection{Future Directions}
Guided by our philosophy of practical applicability, we envision several compelling avenues for future work:
\begin{itemize}
    \item \textbf{Generalization to Larger Backbones and Generative Architectures:} While we validated EdgeMerge on the CLIP ViT-B/32 backbone, our scale-normalized delta activation routing framework is mathematically agnostic to network depth, width, or modality. In future work, we plan to evaluate EdgeMerge on larger vision backbones (e.g., CLIP ViT-L/14) and expand it to Large Language Models (LLMs) like LLaMA and Vision-Language models like LLaVA. In LLMs, inter-task weight conflicts are highly concentrated in the high-dimensional Feed-Forward Network (FFN) layers (representing over 60\% of model parameters). We will investigate localizing our Channel-Wise Softmax Gating (CWSG) to the intermediate projection layers of FFNs (e.g., the gate/up projections in SwiGLU layers), which function as high-leverage semantic bottleneck junctions. This will enable training-free, forward-only on-device parameter composition across diverse text generation and reasoning tasks.

    To concretely illustrate this modality generalizability, we outline a highly feasible proof-of-concept experimental design for text-generative LLMs (e.g., LLaMA-3 8B). Consider the problem of merging a sentiment analysis expert (fine-tuned on IMDb) and a translation expert (fine-tuned on WMT). The intermediate down-projection ($W_{down} \in \mathbb{R}^{d_{model} \times d_{ffn}}$) or gate-projection ($W_{gate} \in \mathbb{R}^{d_{ffn} \times d_{model}}$) layers of SwiGLU FFNs are chosen as our strategic choke-point routing junctions. First, we feed a single batch of 32 unlabeled text prompts from each domain through the pre-trained base model to extract hidden features $X_k$ right before the FFN layer. Second, we compute the channel-wise salience vector $S_k$ based on the activation delta $\Delta H_k = X_k (W^{expert} - W^{base})^T$. Third, we normalize and softmax-gated these vectors across tasks to get channel-specific coefficients $\alpha_k$. Finally, we apply Decoupled Scale Routing (DSR) to compose the weights: keeping standard layers merged at a baseline scale $\lambda_{static}=0.20$, and scaling the gated FFN projection layers with a larger decoupled projection scale $\lambda_{proj} = 1.60$ to offset softmax scaling dampening. This exact mathematical translation of EdgeMerge to the language domain represents a highly promising, zero-gradient, zero-inference-overhead pathway for generative model composition.
    \item \textbf{Data-Free Calibration Synthesizers:} While EdgeMerge is exceptionally lightweight, it still requires a tiny set of calibration inputs. We aim to explore completely data-free/synthetic calibration methods (e.g., using synthetic text descriptions or latent-space noise inversion) to compute channel-wise salience without requiring any physical data, further enhancing privacy and deployment simplicity.
    \item \textbf{On-Chip Hardware Co-Design:} We are interested in implementing EdgeMerge directly on specialized microcontrollers and edge hardware (such as Apple Neural Engines or Raspberry Pi TPU hats). Co-designing our forward-only closed-form mathematical equations with on-chip tensor accelerators will enable edge devices to autonomously and dynamically re-merge their own expert parameters on-the-fly under shifting environmental demands.
\end{itemize}

Ultimately, EdgeMerge demonstrates that we do not need complex, fragile, and resource-heavy optimization loops to achieve effective, localized weight routing. Closed-form, forward-only activation analysis represents a highly scalable, robust, and deployable paradigm for weight-space engineering in production environments.
